\begin{textsample}{POS Dim 2 – vem\_ed\_26 – Score 8.00 – vem\_ed\_26/t140.txt}  \label{ex:f2_pos_023}
Osvaldo Succi Junior na diretoria da COIL Connect Em outubro de 2024, Osvaldo Succi Junior, coordenador dos Projetos Colaborativos Internacionais (PCIs / Cesu), passou a integrar a diretoria executiva da COIL Connect.
Trata-se de uma entidade sem fins lucrativos que visa auxiliar as instituições de ensino superior a desenvolverem a Internacionalização em Casa, por meio dos projetos COIL (Collaborative Online International Learning).
O espectro de participantes da COIL Connect é amplo, abrangendo desde professores individuais em \textbf{busca} de parceiros até instituições como o Centro Paula Souza, que está entre as maiores do mundo em Intercâmbios Virtuais, desenvolvendo aproximadamente 120 projetos ao ano nas Fatecs.
Desde 2013, mais de 11.500 fatecanos participaram de aproximadamente 600 PCIs / Cesu.
O reconhecimento internacional pelos números expressivos das Fatecs, pela diversidade de projetos desenvolvidos e pela inserção nas maiores conferências sobre Intercâmbios Virtuais chamou a \textbf{atenção} do criador da abordagem COIL, Jon Rubin, presidente da COIL Virtual Exchange Foundation (Fundação de Intercâmbios Virtuais do tipo COIL).
Isso resultou em um convite ao Coordenador dos PCIs / Cesu para integrar a diretoria executiva da COIL Connect. “O desafio é imenso e precisa de perspectivas diversas que possibilitem entender como a COIL Connect pode melhor desenvolver sua missão”, comenta Osvaldo Succi Junior. “A abordagem COIL busca auxiliar professores de diferentes países a realizarem projetos, normalmente inseridos em suas disciplinas, em colaboração com seus pares internacionais.
Essa metodologia ajuda os alunos a desenvolverem habilidades de comunicação, letramento digital e trabalho em equipes internacionais, sempre com uma perspectiva simultaneamente local e global”, esclarece Succi Junior.
Para os professores, essa abordagem auxilia na criação de vínculos com instituições internacionais de ensino superior e proporciona oportunidades de participação em \textbf{congressos} e eventos acadêmicos internacionais, escrita de trabalhos em conjunto e desenvolvimento de projetos. continuação Em 2012, Rafael Ferreira Alves (o coordenador técnico da Cesu, na época, era diretor da Fatec Americana) e Osvaldo Succi Junior (então, professor de inglês na unidade) visitaram diversos campi da State University of New York (SUNY) nos Estados Unidos.
Durante esses contatos, foram convidados a participar de um projeto COIL com a SUNY Ulster.
Em 2013, o CPS se tornou a primeira instituição brasileira a utilizar essa abordagem, na Fatec Americana.
Para \textbf{facilitar} o entendimento do conceito COIL, a iniciativa foi rebatizada nas Fatecs como Projetos Colaborativos Internacionais (PCIs / Cesu).
Devido à capilaridade das Fatecs, ao interesse dos professores em interagir com colegas internacionais nos PCIs e à importante \textbf{presença} dos professores de inglês e espanhol na matriz curricular dos cursos, ocorreu uma enorme expansão e hoje as Fatecs do CPS são reconhecidas como uma das maiores instituições do mundo a utilizar essa abordagem. “Com a experiência adquirida, passamos a auxiliar outras instituições internacionais, como a DUOC UC no Chile, a Uniminuto na Colômbia e o Perimeter College nos Estados Unidos, na capacitação de seus professores para o desenvolvimento de PCIs / Cesu.
De simples reprodutores de uma abordagem criada dentro de uma universidade americana, passamos a \textbf{contribuir} para o crescimento e a diversidade dessa abordagem”, ressalta Succi Junior. \textbf{Ficou} claro que cada país e cada universidade têm suas particularidades, e é importante reconhecer, respeitar e ser flexível na ação com colegas internacionais.
Além disso, as vozes do Sul Global, países com maiores dificuldades de recursos, precisam ser ouvidas.
Hoje, a América do Sul tem um papel de destaque no cenário mundial.

% matched lemmas: atenção, busca, congresso, contribuir, facilitar, ficar, presença
\end{textsample}
